
\chapter{View I: the Borcherds--Gritsenko--Nikulin section}
\label{ch:view1-automorphic}

\section{The pentadic level and the Beilinson cut}
\label{sec:view1-pentadic-level}

The automorphic reading of \(\Delta_5\) is the first theorem-level
reading.  It is a holomorphic section of the
weight-\(5\) automorphic line on the genus-two Siegel modular space
\(\Sp_4(\Z)\backslash\mathbb H_2\): the Borcherds--Gritsenko--Nikulin lift of the
weight-zero index-one weak Jacobi form \(\phi_{0,1}\), normalized at the
isotropic cusp \(\mathfrak c_\infty\) by the theta-constant leading
coefficient \(64\).  The \(\Sp_4(\Z)\)-equivariant data of View~I
are the multiplier system \(\nu_{\Delta_5}\), the Borcherds
chamber, the divisor support, and the weight.

\paragraph{Three names.}
The handle ``\(\Delta_5\)'' is shared among three structurally distinct
objects and must be tagged on first use:
\begin{itemize}
\item \emph{the Borcherds--Gritsenko--Nikulin section}
\(\mathcal D_X=\Delta_5\) of \(\mathcal L^{5}\otimes\nu_{\Delta_5}\) on
\(\Sp_4(\Z)\backslash\mathbb H_2\), View~I;
\item \emph{the BKM denominator}
\(\operatorname{den}(\mathfrak g_{\Delta_5})=64^{-1}\Delta_5(2Z)\), View~II,
the Borcherds--Kac--Moody denominator algebra;
\item \emph{the protected Pfaffian}
\(\operatorname{Pf}_{\mathrm{prot}}(\mathfrak D_X^{\mathrm{DI}})\),
View~III, constructed in the Dirac--Igusa chapter as a section on the
Pfaffian line of the cosection-reduced compact Hall datum after the
retained D0-HN datum, retained quotient-orientation datum, chosen Weyl
lifts, retained \((O2_{\mathrm{atlas}})\) wall datum, and finite
geometric Pfaffian datum \((P_{\mathrm{fin}})\) of
Appendix~\ref{app:obstruction-discharges}.
\end{itemize}
The bare label \(\Delta_5\) never substitutes for any of the three: it
names the automorphic section unless a reading is specified, while
\(\mathcal D_X\), \(\operatorname{den}(\mathfrak g_{\Delta_5})\), and
\(\operatorname{Pf}_{\mathrm{prot}}(\mathfrak D_X^{\mathrm{DI}})\) live at
distinct categorical-dimensional levels.

\paragraph{Beilinson cut: scope of View~I.}
View~I is the scalar modular shadow.  In the reconstitution order
``primitive objects first, cross-level identities second, scalar
shadows last'' View~I is a single holomorphic section of an automorphic
line on a Shimura variety, with its multiplier system and its weight.
No operator-level datum is asserted at View~I.  In particular, the equality
\(\mathcal D_X=\operatorname{Pf}_{\mathrm{prot}}(\mathfrak D_X^{\mathrm{DI}})\) of
View~I and View~III is the Pfaffian--Dirac theorem of
Chapter~\ref{ch:pfaffian-dirac}, relative to the retained data recorded
in Appendix~\ref{app:obstruction-discharges}; it is not asserted in this
chapter.
The shadows are \(\Delta_5\), \(\chi_{10}\), the level-\(n\) BPS partition function
\(Z_{\mathrm{BPS}}^{K3\times E}=\Delta_5^{-2}\), and the
Oberdieck--Pixton scalar branch
\(Z^X_{\mathrm{OP}}=-4096\,\Delta_5^{-2}\); none of them is a Hilbert
space, a Hall pairing, an orientation line, or an operator product.

\section{The genus-two Siegel space and the Borcherds chamber}
\label{sec:view1-siegel-and-chamber}

Let \(\mathbb H_2\) be the genus-two Siegel upper half space.  An element
\(Z\in\mathbb H_2\) is a complex symmetric \(2\times 2\) matrix with
positive-definite imaginary part,
\[
Z=\begin{pmatrix}z_1&z_2\\z_2&z_3\end{pmatrix},
\qquad
\Im Z>0.
\]
Set Fourier variables
\[
q=e^{2\pi i z_1},\qquad
r=e^{2\pi i z_2},\qquad
s=e^{2\pi i z_3}.
\]
The fixed cusp polarization at the rational isotropic boundary
\(\mathfrak c_\infty\) of \(\Sp_4(\Z)\backslash\mathbb H_2\) — the standard
type-II cusp polarization — is the
ordered limit
\[
0<|s|\ll|q|\ll|r|^{-1}\ll1.
\]
Equivalently, \(\Im z_3\) goes to \(+\infty\) faster than \(\Im z_1\)
goes to \(+\infty\), and \(\Im z_2\) is bounded.  This polarization
selects a Weyl chamber on the \(O(3,2)\) symmetric domain through the
\(\PSp_4\simeq\SO(3,2)_+\) bridge of Chapter~\ref{ch:singular-theta-lift},
and hence selects a positive cone, the effective semigroup
\[
\Gamma_{\mathrm{eff}}
=\{(n,l,m)\mid m>0,\ n\ge0,\ l\in\Z\}
\cup\{(n,l,0)\mid n>0,\ l\in\Z\}
\cup\{(0,l,0)\mid l<0\}
\subset\Z^3,
\]
and a corresponding completed Fourier expansion ring.  In this completion every
Fourier coefficient of every formal product receives only finitely many
contributions: the chamber walls are the inequalities defining the
ordering, and the chamber-completed product is the only product
considered.

\paragraph{Choices fixed once and for all.}
The cusp \(\mathfrak c_\infty\), the Fourier variables \((q,r,s)\), the
polarization \((0<|s|\ll|q|\ll|r|^{-1}\ll1)\), and the effective semigroup
\(\Gamma_{\mathrm{eff}}\subset\Z^3\) are fixed.  Every View-I product expansion
is the product expansion in the cusp-completed
semigroup algebra of \(\Gamma_{\mathrm{eff}}\).  The cusp expansion at any
other rational isotropic cusp is obtained by an
\(\Sp_4(\Z)\)-translate of these data.

\section{The weak Jacobi input and the K3 elliptic genus}
\label{sec:view1-jacobi-input}

The Borcherds input is the \emph{weight-zero index-one weak Jacobi form}
\(\phi_{0,1}\in J^{\mathrm{weak}}_{0,1}\) of Eichler--Zagier
\cite[\S2.2]{EZ:1}, normalized so that the K3 elliptic genus is
\[
Z_{K3}(\tau,z)=2\,\phi_{0,1}(\tau,z).
\]
Write the Fourier expansion as
\[
\phi_{0,1}(\tau,z)=\sum_{n\ge 0,\;l\in\Z}f(n,l)\,q^nr^l,
\qquad q=e^{2\pi i\tau},\quad r=e^{2\pi iz}.
\]
The first Fourier coefficients are
\begin{equation}
\label{eq:view1-phi01-coefficients}
\begin{aligned}
f(0,-1)&=1,&\quad f(0,0)&=10,&\quad f(0,1)&=1,\\
f(1,-2)&=10,&\quad f(1,-1)&=-64,&\quad f(1,0)&=108,&\quad f(1,1)&=-64,&\quad f(1,2)&=10,\\
f(2,-3)&=1,&\quad f(2,-2)&=108,&\quad f(2,-1)&=-513,&\quad f(2,0)&=808,
\end{aligned}
\end{equation}
extending in the standard way by hyperbolic symmetry
\(f(n,l)=f(n,-l)\) and Eichler--Zagier support
\(f(n,l)=0\) for \(4n-l^2<-1\) \cite[\S2.2]{EZ:1}.  The values in
\eqref{eq:view1-phi01-coefficients} are obtained by the exact rational
expansion of the standard theta-quotient expression
\(\phi_{0,1}=4\sum_{i=2,3,4}\theta_i(\tau,z)^2/\theta_i(\tau,0)^2\)
[\textit{computed}; cross-checked against \cite[Theorem~9.2]{EZ:1}].

\paragraph{Borcherds-weight handle.}
For a fixed Borcherds input \(\Phi\), the weight identity
\(\kappa_{\mathrm{BKM}}(\Phi)=c_{\Phi}(0)/2\) (Borcherds 1995
\cite[Theorem~10.1]{Borcherds95}; Gritsenko 2018
\cite{Gritsenko2018Reflective}) gives, on the row \((t,N)=(1,1)\) of
the Cl\'ery--Gritsenko catalogue \cite[Theorem~1.4]{GC},
\[
c_1(0):=f(0,0)=10,\qquad
\kappa_{\mathrm{BKM}}(\Delta_5)=\frac{c_1(0)}{2}=5,
\]
to be matched in
Proposition~\ref{prop:view1-borcherds-gritsenko-weight-identity} below
with the geometric weight of the Borcherds lift.  The notation
\(\kappa_{\mathrm{BKM}}\) is the BKM cusp-form-weight entry of the
\(\kappa\)-tuple \((\kappa_{\mathrm{cat}},\kappa^{\mathrm{Hodge}}_{\mathrm{ch}},
\kappa^{\mathrm{Heis}}_{\mathrm{ch}},\kappa_{\mathrm{BKM}},\kappa_{\mathrm{fiber}})
=(0,0,3,5,24)\) of the \(K3\times E\) target.  The bare symbol
\(\kappa\) is not a scalar invariant; every entry carries its semantic subscript;
the additive formula
\(\kappa_{\mathrm{BKM}}=\kappa_{\mathrm{ch}}+\chi(\mathcal O_{\mathrm{fiber}})\)
does not define the BKM coordinate.  The BKM coordinate is
\(\kappa_{\mathrm{BKM}}(\Delta_5)=c_1(0)/2=5\).

\section{The theta-constant normalization}
\label{sec:view1-theta-normalization}

The theta-constant Igusa form \(\Delta_5\) is fixed by its leading
Fourier monomial at \(\mathfrak c_\infty\).  Writing the theta-constant
representation of Igusa \cite[\S5]{Igusa1962}, \cite[\S2]{Igusa1964II}
\[
\Delta_5(Z)
=\!\!\!\!\prod_{\substack{a,b\in(\Z/2\Z)^2\\{}^t\!ab\equiv 0\bmod 2}}\!\!\!
\nu_{a,b}(Z),
\qquad
\nu_{a,b}(Z)=\!\sum_{l\in\Z^2}\exp\!\bigl(\pi i\bigl(z[l+\tfrac12 a]+{}^tbl\bigr)\bigr),
\]
the four characteristics with \(a=(0,0)\) start with the constant \(1\),
the two factors with \(a=(1,0)\) start with \(2q^{1/8}\), the two factors
with \(a=(0,1)\) start with \(2s^{1/8}\), and the two factors with
\(a=(1,1)\) start with \(2q^{1/8}r^{1/4}s^{1/8}\); multiplying gives

\begin{proposition}[Theta-constant leading coefficient]
\label{prop:view1-theta-constant-leading-coefficient}
\[
[q^{1/2}r^{1/2}s^{1/2}]\Delta_5\;=\;2^6\;=\;64.
\]
\end{proposition}

\noindent
[\textit{proved}; Igusa 1962 \cite[Theorem~3]{Igusa1962}, Gritsenko
1999 \cite[Theorem~4.1]{GN}; the rational theta-product expansion gives
the same leading coefficient.]

\begin{definition}[Monic Borcherds--Gritsenko form]
\label{def:view1-monic-Borcherds-Gritsenko-form}
The \emph{monic genus-two Borcherds--Gritsenko form} is
\[
D_5(Z)\;:=\;64^{-1}\Delta_5(Z).
\]
By Proposition~\ref{prop:view1-theta-constant-leading-coefficient},
\(D_5=q^{1/2}r^{1/2}s^{1/2}+\ldots\) at \(\mathfrak c_\infty\).
\end{definition}

\noindent
This is the normalization used by Gritsenko--Nikulin
\cite[Theorem~2.1, (2.15)--(2.16)]{GNII}: their Borcherds product is the
monic form \(D_5\); the theta-constant convention restores the prefactor
\(64\).

\section{The Borcherds--Gritsenko--Nikulin section}
\label{sec:view1-borcherds-gritsenko-nikulin-section}

\subsection{The product expansion}
\label{ssec:view1-product-expansion}

Borcherds's automorphic product theorem
\cite[Theorem~10.1]{Borcherds95}, specialized to the genus-two Siegel
modular variety \(\Sp_4(\Z)\backslash\mathbb H_2\) by Gritsenko--Nikulin
\cite[Theorem~2.1, Example~2.4]{GNII} via the
exceptional exterior-square isogeny \(\PSp_4(\R)\simeq\SO(3,2)_+\) of
Chapter~\ref{ch:singular-theta-lift}, lifts the weak Jacobi form
\(\phi_{0,1}\) to a meromorphic Siegel modular form of weight
\(f(0,0)/2=5\) and multiplier system \(\nu_{\Delta_5}\) with infinite
product expansion in the cusp chamber.

\begin{theorem}[Borcherds--Gritsenko--Nikulin lift, monic form]
\label{thm:view1-borcherds-gritsenko-nikulin-lift}
In the cusp chamber \(0<|s|\ll|q|\ll|r|^{-1}\ll1\),
\[
D_5(Z)
\;=\;q^{1/2}r^{1/2}s^{1/2}\!\!
\prod_{(n,l,m)\in\Gamma_{\mathrm{eff}}}\!\!
(1-q^nr^ls^m)^{f(nm,l)},
\]
the product converging absolutely on a neighborhood of \(\mathfrak
c_\infty\) and continuing meromorphically to all of
\(\mathbb H_2\) as a holomorphic Siegel modular form of weight \(5\) on
\(\Sp_4(\Z)\) with character \(\nu_{\Delta_5}\) of order two.
\end{theorem}

\noindent
[\textit{proved}; Borcherds 1995 \cite[Theorem~10.1]{Borcherds95},
Gritsenko--Nikulin 1998 \cite[Theorem~2.1, (2.15)--(2.16)]{GNII},
Gritsenko 1999 \cite[Theorem~4.1]{GN}.]

\begin{corollary}[Theta-normalized section]
\label{cor:view1-theta-normalized-section}
\[
\Delta_5(Z)
\;=\;64\,q^{1/2}r^{1/2}s^{1/2}\!\!
\prod_{(n,l,m)\in\Gamma_{\mathrm{eff}}}\!\!
(1-q^nr^ls^m)^{f(nm,l)}.
\]
\end{corollary}

\noindent
The exponent \(f(nm,l)\) of the \((n,l,m)\)-factor depends only on the
Hecke--Verlinde--DMVV product \(nm\) and on the Jacobi index \(l\),
reflecting the second-quantized structure of the Borcherds lift, whose
combinatorial origin is the Dijkgraaf--Moore--Verlinde--Verlinde
elliptic genus of symmetric products \cite{DMVV}.

\subsection{The View~I datum}
\label{ssec:view1-package}

\begin{definition}[The Borcherds--Gritsenko--Nikulin section]
\label{def:view1-borcherds-gritsenko-nikulin-section}
The \emph{Borcherds--Gritsenko--Nikulin section} of the weight-\(5\)
automorphic line is the View~I representative
\[
\mathcal D_X\;:=\;\Delta_5
\]
on \(\Sp_4(\Z)\backslash\mathbb H_2\), with Borcherds product expansion
at \(\mathfrak c_\infty\) given by Corollary~\ref{cor:view1-theta-normalized-section}.
\end{definition}

\begin{remark}[\(\mathcal D_X\) is the View~I name]
\label{rem:view1-D-X-naming}
The handle \(\mathcal D_X\) is the View~I name of the section.  It coincides as
a holomorphic Siegel modular form with the theta-constant Igusa cusp
form \(\Delta_5\); the renaming records the role: View~I scalar
shadow, the determinant of the virtual K3 index in
\(K_0(\mathrm{sVect}_{\mathrm{fd}})\) (the
\(K_0\)-determinant identity proved at
Proposition~\ref{prop:bgn-identification-recap}; the operator
Pfaffian datum belongs to View~III), and the input section of the Pfaffian--Dirac
theorem of Chapter~\ref{ch:pfaffian-dirac}.  The compact
operator-level datum \(\operatorname{Pf}_{\mathrm{prot}}(\mathfrak
D_X^{\mathrm{DI}})\) of View~III is asserted equal to \(\mathcal D_X\) under
\((D_0)\), \((O_1)\), \((O_1)^+\), \((O_2)\); without these obstruction
trivializations, only \(\mathcal D_X\) is rigorously available.
\end{remark}

\section{Weight, multiplier, and divisor}
\label{sec:view1-weight-multiplier-divisor}

\subsection{The Borcherds weight}
\label{ssec:view1-borcherds-weight}

\begin{proposition}[Borcherds--Gritsenko weight identity at \(N=1\)]
\label{prop:view1-borcherds-gritsenko-weight-identity}
\[
\operatorname{wt}(\mathcal D_X)
\;=\;\frac{f(0,0)}{2}
\;=\;5
\;=\;\frac{c_1(0)}{2}
\;=\;\kappa_{\mathrm{BKM}}(\Delta_5).
\]
The weight is integral; no metaplectic cover enters the weight factor.
\end{proposition}

\noindent
[\textit{proved}; Borcherds 1995 \cite[Theorem~10.1, weight
formula]{Borcherds95}; Gritsenko--Nikulin 1998
\cite[Example~2.4]{GNII}; Gritsenko 1999 \cite[Theorem~4.1, weight]{GN};
the Borcherds weight identity \(\kappa_{\mathrm{BKM}}(\Phi)=c_{\Phi}(0)/2\)
for a fixed input form and lattice; Cl\'ery--Gritsenko \cite{GC}, Vol~III.]

\noindent
The row \((t,N;\operatorname{wt},c)=(1,1;5,10)\) is the leading
\(F_j\)-entry of the eight-row Cl\'ery--Gritsenko catalogue
\cite[Theorem~1.4]{GC}.  In the \(F_j\)-order the remaining rows carry
weights \((2,1,\tfrac12,3,2,\tfrac32,1)\) and constants
\((4,2,1,6,4,3,2)\); in the input-pair order
\((2,1),(1,2),(3,1),(1,3),(4,1),(1,4),(2,2)\) the same data read
\((2,3,1,2,\tfrac12,\tfrac32,1)\) and
\((4,6,2,4,1,3,2)\).  Each row satisfies
\(\kappa_{\mathrm{BKM}}(F_j)=c_j(0)/2\).  The half-integer rows are
realized by \(v_\eta^3\times v_H\) and
\(\chi_4^{(4)}\times v_H\), respectively.  The Vol~III
universal Borcherds-weight identity replaces the additive formula
\(\kappa_{\mathrm{BKM}}=\kappa_{\mathrm{ch}}+\chi(\mathcal O_{\mathrm{fiber}})\):
the \(\kappa\)-tuple records separate invariants, and the
\(K3\times E\) entry is precisely \(\kappa_{\mathrm{BKM}}=5\).

\subsection{The Maass multiplier system}
\label{ssec:view1-maass-multiplier}

The character \(\nu_{\Delta_5}\colon\Sp_4(\Z)\to\C^\times\) of Maass
\cite{MAASS:1} has order two,
\(\nu_{\Delta_5}^{2}=1\), and is determined on the standard generators
of \(\Sp_4(\Z)\) by
\[
\nu_{\Delta_5}\!
\begin{pmatrix}0&\Id_2\\-\Id_2&0\end{pmatrix}=1,
\quad
\nu_{\Delta_5}\!
\begin{pmatrix}\Id_2&B\\0&\Id_2\end{pmatrix}=(-1)^{b_1+b_2+b_3},
\]
\[
\nu_{\Delta_5}\!
\begin{pmatrix}{}^tA^{-1}&0\\0&A\end{pmatrix}
=(-1)^{(1+a_1+a_4)(1+a_2+a_3)+a_1a_4},
\]
with
\(B=\bigl(\begin{smallmatrix}b_1&b_2\\b_2&b_3\end{smallmatrix}\bigr)\)
and \(A=\bigl(\begin{smallmatrix}a_1&a_2\\a_3&a_4\end{smallmatrix}\bigr)\)
\cite{MAASS:1}.  The square \(\Delta_5^2=\Delta_{10}\) is scalar.  The
multiplier records the obstruction to identifying \(\Delta_5\) with an
ordinary modular form on \(\Sp_4(\Z)\): \(\Delta_5\) is a section of
\(\mathcal L^5\otimes\nu_{\Delta_5}\), where \(\mathcal L\) is the Hodge
automorphic line and \(\nu_{\Delta_5}\) is the order-two character; the
section \(\Delta_5\) does not descend to a section of the bare Hodge
line \(\mathcal L^5\), but its square \(\Delta_5^2\) is a scalar
section of \(\mathcal L^{10}\).
The finite automorphic packet
\[
\texttt{certificates/automorphic/delta5\_maass\_character}
\]
records the six chamber generator values, the relations
\(\nu_{\Delta_5}|_{W^{(2)}}=\det\),
\(\nu_{\Delta_5}|_{\Aut(\Poly_{II})}=1\), and
\(\nu_{\Delta_5}^2=1\), and the line row
\(\Delta_5\in H^0(\mathbb H_2,\mathcal L^5\otimes\nu_{\Delta_5})\).
The same packet excludes using \(\nu_{\Delta_5}\) as a Pfaffian
orientation line or as an OP scalar branch.

\subsection{The diagonal divisor}
\label{ssec:view1-diagonal-divisor}

Let
\[
\mathcal H_{\mathrm{diag}}
=\Bigl\{\bigl(\begin{smallmatrix}a&0\\0&b\end{smallmatrix}\bigr)\;\Big|\;
a,b\in\mathbb H_1\Bigr\}
\subset\mathbb H_2
\]
be the locus of diagonal genus-two periods.  The Cl\'ery--Gritsenko
classification \cite{GC} records
\[
\operatorname{supp}\operatorname{div}(\Delta_5)
=\Sp_4(\Z)\cdot\mathcal H_{\mathrm{diag}},\qquad
\operatorname{mult}_{\mathcal H_{\mathrm{diag}}}(\Delta_5)=1,
\]
so the divisor of \(\Delta_5\) is the simple paramodular Humbert
divisor.  In the cusp-chamber language, the type-II reflective walls
\(\mathcal H_{\delta_i}\), \(i=1,2,3\), of the lattice
\(\La^{2,1}_{II}\) (the type-II sublattice of the Borcherds even
lattice \(\La^{2,1}\)) are the components of
\(\Sp_4(\Z)\cdot\mathcal H_{\mathrm{diag}}\) that meet the chosen
fundamental Weyl chamber \(\Poly_{II}\) (Chapter~\ref{ch:view2-bkm}); the
corresponding simple roots
\(\delta_1=2f_2-f_3\), \(\delta_2=2f_{-2}-f_3\), \(\delta_3=f_3\)
have Gram pairings
\((\delta_i,\delta_i)=2\) and \((\delta_i,\delta_j)=-2\) for \(i\ne j\)
\cite[Example~2.4]{GNII}, and the corresponding Borcherds divisor
orders are
\[
d_{\delta_1}^{\mathrm{aut}}=f(0,1)=1,\qquad
d_{\delta_2}^{\mathrm{aut}}=f(0,1)=1,\qquad
d_{\delta_3}^{\mathrm{aut}}=f(0,-1)=1.
\]
[\textit{proved}; Cl\'ery--Gritsenko 2011 \cite[Theorem~1.4]{GC};
Gritsenko--Nikulin 1998 \cite[Example~2.4]{GNII}; the Gram computation
is the one displayed in Appendix~\ref{app:lattice-II21}.]
The finite automorphic divisor packet
\[
\texttt{certificates/automorphic/delta5\_humbert\_divisor}
\]
records the three type-II support representatives, the simple order-one
rows for \(\Delta_5\), the pole-order-two row for \(\Delta_5^{-2}\),
and a separate support-versus-multiplicity row.  It is not a Pfaffian
wall chart, an \(O2\) atlas, or a mirror-period discriminant.

\subsection{Maass-character values on the type-II walls}
\label{ssec:view1-maass-character-walls}

\begin{lemma}[Maass-character values on the type-II reflections]
\label{lem:view1-maass-multiplier-restriction}
With the simple type-II roots
\(\delta_1=2f_2-f_3\), \(\delta_2=2f_{-2}-f_3\), \(\delta_3=f_3\)
and chamber-permuting type-I roots
\(\alpha_1=f_{-2}-f_2\), \(\alpha_2=f_2-f_3\), \(\alpha_3=f_{-2}-f_3\),
\[
\nu_{\Delta_5}(s_{\delta_i})\;=\;-1,\qquad
\nu_{\Delta_5}(s_{\alpha_i})\;=\;+1,\qquad i=1,2,3.
\]
Equivalently, for chosen \(\Sp_4\)-lifts and the weight-five slash
action without character,
\[
(\Delta_5|_5s_{\delta_i})(Z)=-\Delta_5(Z),\qquad
(\Delta_5|_5s_{\alpha_i})(Z)=+\Delta_5(Z).
\]
\end{lemma}

\noindent
[\textit{proved}; Maass \cite{MAASS:1}; Gritsenko--Nikulin
\cite[\S3]{GN}; Gritsenko \cite{Gritsenko99}.  The type-II values
follow from \(d_{\delta_i}^{\mathrm{aut}}=1\) and the chamber-wall
cocycle calculation of
Lemma~\ref{lem:bkm-maass-restriction}, which fixes
\(\nu_{\Delta_5}|_{W^{(2)}(\La^{2,1}_{II})}=\det\) on the three
type-II generators; the type-I values follow from
\(d_{\alpha_i}^{\mathrm{aut}}=0\) (vanishing of \(f(0,\pm 2)\) by the
support of \(\phi_{0,1}\) at \(q=0\)).]

The chamber-permuting reflections
\(s_{f_{-2}-f_2},s_{f_2-f_3},s_{f_{-2}-f_3}\) act trivially on
\(\Delta_5\) because the corresponding type-I roots
\(\alpha_i\) are not \(\Sp_4\)-images of timelike short roots in the
cusp chamber: type-I square-$2$ vectors are outside the type-II
sublattice $\La^{2,1}_{II}$ on which the Borcherds product
expansion of $\Delta_5$ is supported, so their walls do not appear
in $\operatorname{div}(\Delta_5)$ at all
\cite[\S2.2]{EZ:1}.

\section{Identification with the virtual \(K_0\)-determinant}
\label{sec:view1-K0-determinant-identification}

The Fock superdeterminant of the Borcherds-graded super vector space
\(\mathcal V=\bigoplus_{(n,l,m)\in\Gamma_{\mathrm{eff}}}
\mathbb U_{nm,l}\) — the \(K_0\)-class of the second-quantized
Hecke--DMVV half of the K3 elliptic genus, with the cusp/Weyl boundary
factors at \(m=0\) carrying the input multiplicities \(f(0,l)\) and
the bulk factors at \(m>0\) carrying \(f(nm,l)\) on
\(\Gamma_{\mathrm{eff}}\) (Definition~\ref{def:gram-index-lattice}) — is

\begin{lemma}[\(K_0\)-Fock superdeterminant]
\label{lem:view1-K0-Fock-superdeterminant}
\[
\sdet_{\mathcal V}(1-x)
\;=\;\!\!\prod_{(n,l,m)\in\Gamma_{\mathrm{eff}}}\!\!
(1-q^nr^ls^m)^{f(nm,l)},
\qquad x^{(n,l,m)}=q^nr^ls^m.
\]
\end{lemma}

\noindent
[\textit{proved}; the Fock-trace lemma is the three-way comparison
\(\sdet_\mathcal V(1-x)=\prod(1-x^\gamma)^{\sdim\mathcal V_\gamma}\)
under the Borcherds-graded second-quantization; the
automorphic--orthogonal--BKM identification sits in
Theorem~\ref{thm:bgn-three-way} of
Chapter~\ref{ch:singular-theta-lift}.]

\begin{proposition}[BGN identification of the virtual determinant]
\label{prop:bgn-identification-recap}
With \(v_0(Z):=64\,q^{1/2}r^{1/2}s^{1/2}\) the cusp/Weyl monomial of
Chapter~\ref{ch:k3xe-setup}, the normalized \(K_0\)-determinant agrees with
the View~I section,
\[
\mathcal D_X(Z)
\;:=\;v_0(Z)\,\sdet_{\mathcal V}(1-x)
\;=\;\Delta_5(Z),
\]
in the cusp-chamber expansion at \(\mathfrak c_\infty\) and hence as
sections of \(\mathcal L^5\otimes\nu_{\Delta_5}\).
\end{proposition}

\begin{proof}
Combine
Theorem~\ref{thm:view1-borcherds-gritsenko-nikulin-lift}, the
theta-constant prefactor
Proposition~\ref{prop:view1-theta-constant-leading-coefficient},
Lemma~\ref{lem:view1-K0-Fock-superdeterminant}, and Cohen's
\(q\)-expansion principle for genus-two Siegel modular forms with
half-integral character: a meromorphic Siegel form is determined by its
expansion at one isotropic cusp.  Both sides have the same Borcherds
product expansion at \(\mathfrak c_\infty\) by direct comparison of
exponents, and \(\mathcal D_X\) is holomorphic of weight \(5\) by
Theorem~\ref{thm:view1-borcherds-gritsenko-nikulin-lift}.
\end{proof}

\noindent
[\textit{proved}; Borcherds 1995 \cite[Theorem~10.1]{Borcherds95},
Gritsenko--Nikulin 1998 \cite[Theorem~2.1]{GNII}, Igusa 1962/1964
\cite{Igusa1962,Igusa1964II}.]

\begin{remark}[Scope of the BGN identification]
\label{rem:view1-bgn-scope}
Proposition~\ref{prop:bgn-identification-recap} is an equality of
holomorphic sections of an automorphic line bundle.  The carrier
\(\mathcal V\) is a class in
\(K_0(\Gamma_{\mathrm{eff}}\text{-graded super vector spaces})\), with no
vertex, chiral, or holomorphic \(E_3\)-factorization structure.
The passage from the \(K_0\)-Fock determinant to a section of an
oriented Pfaffian line over the cosection-reduced compact Hall datum,
and the identification of that section with \(\mathcal D_X\), is the
Pfaffian--Dirac theorem of Chapter~\ref{ch:pfaffian-dirac}, relative to
the retained data recorded in Appendix~\ref{app:obstruction-discharges}.
View~I supplies only the automorphic determinant section.
\end{remark}

\section{The squared determinant section and OP normalization}
\label{sec:view1-squared-scalar}

The Maass character has order two, so the square
\(\Delta_5^2=\Delta_{10}\) is a scalar weight-\(10\) Igusa form on
\(\Sp_4(\Z)\), with no multiplier system.

\begin{proposition}[The square is the Oberdieck--Pandharipande $\chi_{10}$]
\label{prop:view1-square-chi-10}
\[
\Delta_5^2\;=\;\Delta_{10}\;=\;\Phi_{10}^{\mathrm{un}}\;=\;\chi_{10}.
\]
\end{proposition}

\noindent
[\textit{proved}; Igusa \cite{Igusa1962,Igusa1964II}; the symbol
\(\Phi_{10}^{\mathrm{un}}\) is Borcherds's name for the squared theta-constant Igusa
form, \(\chi_{10}\) is the Oberdieck--Pandharipande
\cite{OPand,OB} name; both equal \(\Delta_5^2\).  The naming is
\(\Phi_{10}^{\mathrm{un}}=\chi_{10}=\Delta_5^2\), with the Oberdieck--Pixton OP
normalization \(\chi_{10}^{\mathrm{OP}}=D_5^{\,2}=4096^{-1}\Delta_5^2\)
the monic-on-\(D_5\) variant.]

\paragraph{The OP scalar branch.}
With the OP normalization \(\chi_{10}^{\mathrm{OP}}=D_5^2\), the
Oberdieck--Pixton scalar partition function on \(K3\times E\) is
\cite[(0.4) and Theorem~1]{OB}
\[
Z_{\mathrm{OP}}^{X}\;=\;-\,(\chi_{10}^{\mathrm{OP}})^{-1}
\;=\;-\,4096\,\Delta_5^{-2}.
\]
The factor \(-4096=-(64)^2\) records the squared theta-constant leading
coefficient and the OP scalar sign convention.  At View~I this is the
closed protected trace \(\Delta_5^{-2}\); the orientation character
which distinguishes the Pfaffian section \(\Delta_5\) from its negative
is lost under squaring.  The operator datum entering that trace is
defined at the Pfaffian and orientation level relative to the retained
data of Appendix~\ref{app:obstruction-discharges}.

\paragraph{Three names for the squared form, distinguished.}
The notation is fixed as follows:
\(\Delta_5^2=\Delta_{10}\) is the Igusa convention (theta-constant
square); \(\Phi_{10}^{\mathrm{un}}=\Delta_5^2\) is the
Borcherds--Eisenstein convention;
\(\chi_{10}=\Delta_5^2\) is the Oberdieck--Pandharipande convention; and
\(\chi_{10}^{\mathrm{OP}}=4096^{-1}\Delta_5^2=D_5^{\,2}\) is the
Oberdieck--Pixton variant in which \(D_5\) is the monic Borcherds form
(Definition~\ref{def:view1-monic-Borcherds-Gritsenko-form}).
The finite packet
\(\texttt{certificates/normalizations/delta5\_theta\_leading}\)
records these as separate convention rows and relation rows; equality
of the unscaled sections is not used to identify the OP monic form or
either inverse scalar branch.

\section{The View~II shadow: the \texorpdfstring{\(2Z\)}{2Z} denominator}
\label{sec:view1-2Z-denominator}

The Gritsenko--Nikulin denominator algebra of the
Borcherds--Kac--Moody Lie superalgebra \(\mathfrak g_{\Delta_5}\)
arises as the View~II shadow of \(\Delta_5\) under the doubling
\(Z\mapsto 2Z\):

\begin{proposition}[Gritsenko--Nikulin doubled denominator]
\label{prop:view1-GN-doubled-denominator}
\[
\operatorname{den}(\mathfrak g_{\Delta_5})\;=\;D_5(2Z)\;=\;64^{-1}\Delta_5(2Z).
\]
\end{proposition}

\noindent
[\textit{proved}; Gritsenko--Nikulin 1998 \cite[Theorem~2.1,
(2.15)--(2.16); \S5]{GNII}; Gritsenko 1999 \cite[Theorem~4.1]{GN}.]

\noindent
The doubling records the embedding
\(\La^{2,1}_{II}=\alpha(\Gamma_{\mathrm{ind}})\subset\La^{2,1}\) of
index \(4\).  Its type-II walls are primitive roots of square \(2\), and
the substitution \(Z\mapsto 2Z\) replaces
\(\exp(-\pi i(\rho,Z))\) by \(\exp(-2\pi i(\rho,Z))\) in the
Gritsenko--Nikulin denominator normalization
\cite[Theorem~2.1, (2.15)]{GNII}.  At View~I this is a pointer; the
Cartan, Chevalley, real-and-imaginary roots, and Weyl--Kac--Borcherds
denominator identity of \(\mathfrak g_{\Delta_5}\) sit in
Chapter~\ref{ch:view2-bkm}.

\section{The \(\PSp_4\simeq\SO(3,2)_+\) bridge to View~V}
\label{sec:view1-Sp-Spin-bridge}

The Borcherds construction is naturally written on the orthogonal
side: \(\Delta_5\) is the regularized theta lift of
\(\phi^{\mathrm{Weil}}_{0,1}\) on \(\La^{3,2}\), pulled back to
\(\Sp_4\) through the exterior-square isogeny
\[
\PSp_4(\mathbb R)\;\simeq\;\SO(3,2)_+ .
\]
The corresponding spin cover is the real group with Lie algebra
\(\mathfrak{sp}_4(\mathbb R)\).  At the symbolic level View~I
(\(\Delta_5\) on \(\Sp_4(\Z)\backslash\mathbb H_2\)) and View~V
(\(\Delta_5\) as the singular theta lift on
\(\Gamma_{3,2}\backslash\Hpl^{\IV}_+\)) are the two sides of this bridge: the
weight-\(5\) automorphic line on \(\Sp_4(\Z)\backslash\mathbb H_2\) is
the pullback of the weight-\(5\) orthogonal modular line on
\(\Gamma_{3,2}\backslash\Hpl^{\IV}_+\), and the multiplier \(\nu_{\Delta_5}\)
restricts to the orthogonal Maass character.  The full theta-lift
construction — Borcherds' \emph{singular theta correspondence}
\cite{Borcherds95}, \cite{B} — produces the orthogonal modular form;
the automorphic section in View~I is its pullback to the
\(\Sp_4\)-side of the bridge.

\section{Summary of View~I}
\label{sec:view1-summary}

\begin{theorem}[View~I datum]
\label{thm:view1-package}
The genus-two Igusa cusp form \(\Delta_5\) is the Borcherds--Gritsenko--Nikulin
section
\[
\mathcal D_X\;:=\;\Delta_5
\]
of the weight-\(5\) automorphic line \(\mathcal L^{5}\otimes\nu_{\Delta_5}\)
over \(\Sp_4(\Z)\backslash\mathbb H_2\):

\smallskip
\begin{enumerate}
\item[\textup{(i)}]
At the rational isotropic cusp \(\mathfrak c_\infty\) with cusp polarization
\(0<|s|\ll|q|\ll|r|^{-1}\ll1\), \(\mathcal D_X\) admits the Borcherds product
expansion
\[
\mathcal D_X(Z)
\;=\;64\,q^{1/2}r^{1/2}s^{1/2}\!\!
\prod_{(n,l,m)\in\Gamma_{\mathrm{eff}}}\!\!
(1-q^nr^ls^m)^{f(nm,l)},
\]
with leading theta-constant coefficient
\([q^{1/2}r^{1/2}s^{1/2}]\mathcal D_X=64\), and exponents \(f(nm,l)\) the Fourier
coefficients of the K3 weight-zero index-one weak Jacobi form
\(\phi_{0,1}=\frac12 Z_{K3}\) (Eichler--Zagier \cite[\S2.2]{EZ:1}).
\item[\textup{(ii)}]
\(\mathcal D_X\) has weight
\(\operatorname{wt}(\mathcal D_X)=f(0,0)/2=5=c_1(0)/2=\kappa_{\mathrm{BKM}}(\Delta_5)\),
the leading entry \((t,N;\operatorname{wt},c)=(1,1;5,10)\) of the
Cl\'ery--Gritsenko catalogue \cite[Theorem~1.4]{GC}.  The multiplier is
the Maass character \(\nu_{\Delta_5}\) of order two.
\item[\textup{(iii)}]
The divisor of \(\mathcal D_X\) is the simple paramodular Humbert divisor
\(\Sp_4(\Z)\cdot\mathcal H_{\mathrm{diag}}\); the type-II reflective
walls \(\mathcal H_{\delta_i}\) (\(i=1,2,3\)) carry order
\(d_{\delta_i}^{\mathrm{aut}}=1\), and
\(\nu_{\Delta_5}(s_{\delta_i})=-1=(-1)^{d_{\delta_i}^{\mathrm{aut}}}\),
while the chamber-permuting roots \(\alpha_i\) carry order \(0\) and
\(\nu_{\Delta_5}(s_{\alpha_i})=+1\) (Lemma~\ref{lem:view1-maass-multiplier-restriction}).
\item[\textup{(iv)}]
The square \(\mathcal D_X^{\,2}=\Delta_{10}=\Phi_{10}^{\mathrm{un}}=\chi_{10}\) is a scalar
weight-\(10\) Igusa form, and the Oberdieck--Pixton scalar branch is
\(Z_{\mathrm{OP}}^X=-(\chi_{10}^{\mathrm{OP}})^{-1}=-4096\,\Delta_5^{-2}\).
\item[\textup{(v)}]
The View~II shadow is the doubled denominator
\(\operatorname{den}(\mathfrak g_{\Delta_5})=D_5(2Z)=64^{-1}\Delta_5(2Z)\)
(Proposition~\ref{prop:view1-GN-doubled-denominator}); the View~III
operator-level datum
\(\operatorname{Pf}_{\mathrm{prot}}(\mathfrak D_X^{\mathrm{DI}})\) is the
subject of Chapter~\ref{chap:dirac-igusa-pro-object} and
Appendix~\ref{app:obstruction-discharges}.
\end{enumerate}
\end{theorem}

\noindent
[\textit{proved};
Borcherds 1995 \cite[Theorem~10.1]{Borcherds95},
Borcherds 1998 \cite{B},
Gritsenko--Nikulin 1998 \cite[Theorem~2.1, Example~2.4]{GNII},
Gritsenko 1999 \cite[Theorem~4.1]{GN},
Igusa 1962 \cite[Theorem~3]{Igusa1962},
Igusa 1964 \cite[\S2]{Igusa1964II},
Cl\'ery--Gritsenko 2011 \cite[Theorem~1.4]{GC},
Maass 1964 \cite{MAASS:1},
Eichler--Zagier 1985 \cite[\S2.2]{EZ:1},
Oberdieck--Pixton 2018 \cite[(0.4) and Theorem~1]{OB},
Oberdieck--Pandharipande 2016 \cite{OPand}; the lattice and
theta-product computations are displayed in Appendices~\ref{app:lattice-II21}
and~\ref{app:borcherds-product-expansion}.]

\paragraph{What View~I does not assert.}
View~I does not assert any operator-level datum on \(K3\times E\),
does not construct the Borcherds--Kac--Moody algebra
\(\mathfrak g_{\Delta_5}\) (Chapter~\ref{ch:view2-bkm}), does not
construct the protected first-order operator \(\mathfrak D_X^{\mathrm{DI}}\)
or the Pfaffian line \(\mathscr L_{\mathrm{Pf},X}\)
(Chapter~\ref{chap:dirac-igusa-pro-object}), and does not commit to the
\(K3\times E\) realization of the chosen \(E_3\)-prefactorization datum
\(\mathfrak A_{K3\times E}\) at any level above the scalar shadow.  In
the language of the Beilinson cut, View~I is the bottom of the
reconstitution order: a holomorphic Siegel cusp form, with its weight,
multiplier, divisor, and Borcherds product, and nothing more.
